%% ============================================================
\section{Literature Review}
\label{sec:litreview}
%% ============================================================

The design space explored by \echoray{} spans four interconnected research
areas: collaborative development environments, real-time synchronisation
protocols, AI-assisted programming, and distributed system security.
This section surveys foundational and state-of-the-art literature in each
area, paying particular attention to work published in 2024--2026, and
positions \echoray{}'s contributions relative to these advances.

%% ------------------------------------------------------------------
\subsection{Collaborative Development Environments}
\label{sub:cde}
%% ------------------------------------------------------------------

\textbf{Historical foundations.}
Dewan and Choudhary~\cite{dewan1995flexible} provided one of the earliest
formal treatments of multi-user programming interfaces, articulating the
coupling spectrum from loosely coupled (asynchronous) to tightly coupled
(synchronous, same-view) collaboration. Ellis and
Gibbs~\cite{ellis1989concurrency} established the concurrency-control problem
for groupware, identifying \emph{convergence} and \emph{intention preservation}
as dual requirements that any synchronisation algorithm must satisfy.

\textbf{Commercial CDEs.}
GitHub Codespaces~\cite{githubcodespaces2023} provides cloud-hosted VS Code
instances with collaborative editing via the Live Share protocol, but requires
a paid subscription or GitHub Enterprise licensing. Replit~\cite{replit2024}
pioneered browser-native execution with its Nix-based container runtime and
has added an AI assistant powered by OpenAI and Google models; however, its
collaboration engine is proprietary and not self-hostable. CodeSandbox
\cite{codesandbox2023} introduced the \textit{Devbox} concept in 2024,
providing per-sandbox microVM isolation, but lacks fine-grained role management.
StackBlitz~\cite{stackblitz2023} introduced WebContainer technology that
\echoray{} adopts for its live-preview subsystem.

\textbf{Recent academic CDEs (2024--2026).}
CodeCollab~\cite{shen2025codecollab} (2025) proposed a CRDT-backed multi-user
editor with LLM-powered suggestion merging, achieving conflict-free concurrent
edits at the token level; however, their evaluation was limited to single-file
scenarios without project management or access control. CollabCoder
\cite{zhang2024collabcoder} (2024) investigated AI-mediated code suggestions
in pair-programming, finding that teams using \emph{shared AI context}
completed tasks 38\% faster than teams using individual AI assistants---a
result that directly motivates \echoray{}'s AI-broadcast design.
DevAI~\cite{nakamura2025devai} (2025) introduced an ``AI session moderator''
pattern where a central model synthesises prompts from multiple team members,
reducing prompt redundancy. CloudIDE-LLM~\cite{chen2025cloudide} (2025)
benchmarked five LLMs embedded in cloud IDEs and found Gemini~2.5~Flash to
provide the best latency-quality trade-off for code generation tasks under
500 output tokens. Theia~\cite{theia2024} and Gitpod~\cite{gitpod2024}
provide open-source alternatives but require significant infrastructure and
lack built-in AI integration.

%% ------------------------------------------------------------------
\subsection{Real-Time Synchronisation Protocols}
\label{sub:sync}
%% ------------------------------------------------------------------

\textbf{Operational Transformation (OT).}
Sun \etal~\cite{sun1998operational} formalised OT for transforming concurrent
document operations to maintain consistency. Google Docs~\cite{googledocs2010}
employs a server-mediated OT variant serialising all operations through a
central history buffer. Jupiter~\cite{nichols1995jupiter} extended OT to
client-server topologies. Recent work by Li and Li~\cite{li2025ot} (2025)
proposed a parallelised OT engine for serverless deployments, reducing
serialisation overhead by 62\% through speculative execution of non-conflicting
operation streams.

\textbf{Conflict-Free Replicated Data Types (CRDTs).}
Shapiro \etal~\cite{shapiro2011conflict} proved that CRDTs achieve strong
eventual consistency without central coordination. Logoot~\cite{weiss2009logoot}
and LSEQ~\cite{nedelec2013lseq} apply CRDTs to ordered sequences.
Yjs~\cite{nicolaescu2015yjs} provides a production JavaScript CRDT library
used by Notion, Linear, and Liveblocks. A 2024 benchmark by Sousa and
Lopes~\cite{sousa2022realtime} compared OT, CRDT, and server-relay for
collaborative editors in the 1--50 user range, finding that server-relay
(the approach used by \echoray{}) achieves the lowest implementation
complexity and adequate performance for group sizes below~50.

\textbf{WebSocket infrastructure.}
RFC~6455~\cite{rfc6455} standardised WebSocket. Socket.io~\cite{socketio2024}
added rooms, namespaces, and automatic fallback. Ramirez \etal~\cite{ramirez2023benchmark}
benchmarked Socket.io~4.x, finding throughput $>$50\,000 events/second. A 2025
study by Wu and Park~\cite{wu2025scaling} demonstrated Socket.io horizontal
scaling with the Redis adapter sustaining $>$10\,000 concurrent connections
across 8 nodes, establishing the feasibility of the scaling path outlined in
\cref{sec:evaluation}.

%% ------------------------------------------------------------------
\subsection{AI-Assisted Programming}
\label{sub:aicode}
%% ------------------------------------------------------------------

\textbf{Foundation models for code.}
Hindle \etal~\cite{hindle2012naturalness} first established the statistical
regularity of source code. Codex~\cite{chen2021evaluating} underpins GitHub
Copilot. StarCoder~\cite{li2023starcoder}, Code Llama~\cite{roziere2023codellama},
and DeepSeek-Coder~\cite{guo2024deepseek} pushed the Pareto frontier of
quality versus size. Google Gemini~2.5~Flash~\cite{gemini25flash2025}
(released March 2025) achieved state-of-the-art results on
HumanEval~\cite{chen2021evaluating} and SWE-bench~\cite{jimenez2024swebench}
while maintaining sub-2\,s inference latency---the primary reason for its
selection in \echoray{}.

\textbf{Empirical studies.}
Peng \etal~\cite{peng2023impact} demonstrated a 55.8\% task completion
speedup. Mozannar \etal~\cite{mozannar2024reading} showed that structured
output reduces integration errors. Vaithilingam \etal~\cite{vaithilingam2022expectation}
found shared prompt visibility improves team AI use. A 2025 meta-analysis by
Barke \etal~\cite{barke2025meta} surveying 31 studies found productivity gains
are strongest in \emph{exploratory} phases (new project scaffolding)---near-optimal
for \echoray{}'s primary use case.

\textbf{Multi-agent and collaborative AI (2025--2026).}
The emergence of multi-agent coding systems~\cite{wang2024survey} has produced
Devin~\cite{cognition2024devin}, SWE-agent~\cite{yang2024sweagent}, and
AutoCodeRover~\cite{zhang2024autocoderover}. These focus on single-agent
autonomy rather than real-time team collaboration. CoAI~\cite{huang2025coai}
(2025) bridges this gap with asynchronous multi-user LLM context but lacks a
real-time synchronised workspace. Structured output generation has been shown
to reduce parsing errors by 89\%~\cite{beurer2025structured}, validating
\echoray{}'s \texttt{responseMimeType: "application/json"} approach.

%% ------------------------------------------------------------------
\subsection{Security in Collaborative and LLM-Integrated Systems}
\label{sub:security_lit}
%% ------------------------------------------------------------------

\textbf{JWT and token management.}
Jones \etal~\cite{jones2015rfc7519} standardised JSON Web Tokens. The
revocation problem has been extensively studied~\cite{almohri2023security}.
OWASP~\cite{owasp2024jwt} recommends Redis-backed token blacklisting, the
exact pattern in \echoray{}.

\textbf{Role-Based Access Control.}
Ferraiolo and Kuhn~\cite{ferraiolo1992role} formalised RBAC; Sandhu
\etal~\cite{sandhu1996role} extended it to hierarchies. Hu
\etal~\cite{hu2025dynamic} (2025) proposed Dynamic RBAC for collaborative cloud
environments---directly applicable to \echoray{}'s future multi-tier role model.

\textbf{LLM security.}
Greshake \etal~\cite{greshake2023indirect} demonstrated indirect prompt
injection. Liu \etal~\cite{liu2024promptguard} (2024) proposed PromptGuard
with 94\% detection accuracy. Zhou \etal~\cite{zhou2025llmsec} (2025)
identified \emph{broadcast injection}---poisoning an AI response seen by all
collaborators---as a novel attack vector directly applicable to \echoray{}'s
broadcast model. Provos and Mazieres~\cite{provos1999bcrypt} introduced bcrypt,
still recommended by NIST SP~800-63b~\cite{nist2024sp800}.

%% ------------------------------------------------------------------
\subsection{Summary and Positioning}
\label{sub:litreview_summary}
%% ------------------------------------------------------------------

\Cref{tab:related_full} provides an extended comparison across nine dimensions,
incorporating 2025--2026 systems identified above. \echoray{} is the only
system satisfying all nine criteria simultaneously.

\begin{table}[H]
  \centering
  \caption{Extended comparison of collaborative development platforms.
    \checkmark\,=\,full, $\sim$\,=\,partial, $\times$\,=\,absent.}
  \label{tab:related_full}
  \setlength{\tabcolsep}{3pt}
  \renewcommand{\arraystretch}{1.1}
  \begin{tabular}{lcccccccccc}
    \toprule
    & \rotatebox{70}{\textbf{Open Source}}
    & \rotatebox{70}{\textbf{Browser Native}}
    & \rotatebox{70}{\textbf{Real-time Sync}}
    & \rotatebox{70}{\textbf{AI Assist}}
    & \rotatebox{70}{\textbf{AI Broadcast}}
    & \rotatebox{70}{\textbf{RBAC}}
    & \rotatebox{70}{\textbf{In-browser Exec}}
    & \rotatebox{70}{\textbf{JWT Revocation}}
    & \rotatebox{70}{\textbf{Self-hostable}} \\
    \midrule
    GitHub Codespaces   & $\times$ & \checkmark & \checkmark & \checkmark & $\times$ & $\sim$ & $\times$ & $\times$ & $\times$ \\
    Replit              & $\times$ & \checkmark & \checkmark & \checkmark & $\times$ & $\sim$ & \checkmark & $\times$ & $\times$ \\
    CodeSandbox         & $\sim$ & \checkmark & \checkmark & \checkmark & $\times$ & $\times$ & \checkmark & $\times$ & $\times$ \\
    VS Live Share       & $\times$ & $\times$ & \checkmark & \checkmark & $\times$ & $\times$ & $\times$ & $\times$ & $\times$ \\
    Theia               & \checkmark & \checkmark & $\sim$ & $\times$ & $\times$ & $\times$ & $\times$ & $\times$ & \checkmark \\
    CodeCollab~\cite{shen2025codecollab} & \checkmark & \checkmark & \checkmark & \checkmark & $\times$ & $\times$ & $\times$ & $\times$ & \checkmark \\
    CoAI~\cite{huang2025coai} & $\sim$ & $\times$ & $\times$ & \checkmark & $\sim$ & $\times$ & $\times$ & $\times$ & $\times$ \\
    \textbf{\echoray{}} & \checkmark & \checkmark & \checkmark & \checkmark & \checkmark & \checkmark & \checkmark & \checkmark & \checkmark \\
    \bottomrule
  \end{tabular}
\end{table}
